\documentclass[a4paper, 12pt, twoside, openright]{book}
\usepackage[english]{babel}
\usepackage[T1]{fontenc}
\usepackage[latin1]{inputenc}
\usepackage{fancyhdr}
\usepackage{float}
\usepackage{graphicx}
\usepackage{wrapfig}
%------------------------------ colors
\usepackage[usenames,dvipsnames,table]{xcolor} % use colors on table and more
\definecolor{333}{RGB}{51, 51, 51} % define custom color
%------------------------------ source code
\usepackage{listings}
\lstset{
  basicstyle=\footnotesize\sffamily,
  commentstyle=\itshape\color{gray},
  captionpos=b,
  frame=shadowbox,
  language=HTML,
  rulesepcolor=\color{333},
  tabsize=2
}
%------------------------------ define Abstract environment, missing in the 'book' class
\newenvironment{abstract}{\cleardoublepage \null \vfill \begin{center}\bfseries\abstractname \end{center}}{\vfill\null}
\addto\captionsenglish{\renewcommand*\abstractname{Sommario}} % change Abstract title
%------------------------------ active url
\usepackage{url}
\renewcommand{\UrlFont}{\color{black}\small\ttfamily}
\usepackage[colorlinks=true, linkcolor=black, citecolor=black, urlcolor=black]{hyperref} % active ref
%------------------------------ macros
\newcommand{\sectionname}{Section} % define Section ref
\newcommand{\subsectionname}{Sub-section} % define Sub-section ref
\renewcommand*\arraystretch{1.4} % tables padding


\begin{document}
\frontmatter

\begin{titlepage} %------------------------------ TITLE PAGE
\begin{center}
\vbox to0pt{\vbox to\textheight{\vfill \includegraphics[width=11.5cm]{images/unipd-light.png} \vfill}\vss}

\hspace{0.5cm}
\begin{minipage}{.20\textwidth}
  \includegraphics[height=2.5cm]{images/unipd-bn.png}
\end{minipage}\begin{minipage}{.90\textwidth}
  \begin{table}[H]
  \begin{tabular}{l}
  \scshape{\Large{\bfseries{Universit� degli Studi di Padova}}} \\
  \hline \\
  \scshape{\Large{Dipartimento di Matematica}} \\
  \end{tabular}
  \end{table}
\end{minipage}

\vspace{1cm}
\emph{\Large{Laurea~Magistrale~in~Informatica}} \\
\vspace{0.5cm}
\emph{\Large{Indirizzo~Sistemi}} \\
\vspace{1.5cm}
\scshape{\Large{\bfseries{PathS: un servizio Location-Based per percorsi alternativi}}} \\
\vspace{0.2cm} \linespread{1} \scshape{\large{\bfseries{(Studio ed implementazione del sistema server)}}}
\end{center}

\vfill
\begin{normalsize}
\begin{flushleft}
  \hspace{45pt} \textit{Laureando} \hspace{160pt} \textit{Relatore}\\
  \vspace{5pt}
  \hspace{30pt} \large{\textbf{Tobia Zorzan}} \hspace{70pt} \large{\textbf{Prof. Claudio Palazzi}}\\
  \vspace{10pt}
  \hspace{260pt} \normalsize{\textit{Co-relatore}}\\
  \vspace{5pt}
  \hspace{200pt} \large{\textbf{Prof.ssa Ombretta Gaggi}}
\end{flushleft}
\end{normalsize}

\vfill
\begin{center}
\hspace{-0.2cm}
\line(1, 0){360}

\textsc{Anno Accademico 2014/2015}
\end{center}
\end{titlepage}


%-\cleardoublepage % make left page blank
%-\thispagestyle{empty} %------------------------------ DEDICA

%-\null
%-\vspace{2cm}
%-\begin{flushright}
%-A ...
%-\end{flushright}
%-\vfill

%-\begin{quote}
%-  Quote

%-  \textit{Author}
%-\end{quote}
%-\vfill
%-\null


\begingroup %------------------------------ CONTENTS
  \makeatletter
  \let\ps@plain\ps@empty
  \makeatother
  \tableofcontents
  \clearpage
\endgroup


\begin{abstract} %------------------------------ ABSTRACT
\markboth{}{} % remove header
\thispagestyle{empty}
I recenti sviluppi dei dispositivi \emph{smartphone}, in particolare l'inclusione di sensori sempre pi� evoluti, ha portato alla definizione di un nuovo paradigma denominato \emph{Web\textsuperscript{2}} (\emph{Web Squared}). La diffusione capillare di questi strumenti consente di creare una rete di sensori distribuiti in grado di raccogliere una ingente mole di dati. I dati stessi possono essere quindi elaborati e fornire il supporto a nuovi servizi che non si basano sulle sole informazioni \emph{statiche} ma su un contesto \emph{dinamico} modificato dall'interazione degli utenti. Il progetto \emph{PathS} vuole essere un prototipo che segue questo paradigma. La sua componente \emph{client} raccoglie informazioni aggiuntive sui percorsi pedonali percorsi dagli utenti. La parte \emph{server} elabora le informazioni raccolte e mantiene una base di dati in grado di supportare servizi di interrogazione che adottano criteri aggiuntivi a quelli tradizionali.
Si espone quindi quali sono stati i principi e le valutazioni eseguite nella realizzazione del componente \emph{PathS Server}, i risultati raggiunti e le possibili aree di miglioramento individuate.
\end{abstract}


\mainmatter

\chapter{Introduzione} %------------------------------ INTRODUZIONE
\thispagestyle{empty}

Introduzione al progetto. Presentazione dello scopo principale e degli attori coinvolti. Rassegna del contenuto dei capitoli.

\chapter{Contesto e lavori correlati} %------------------------------ CONTESTO
\thispagestyle{empty}

\section{Web\textsuperscript{2} e Opportunistic Sensing}
Presentazione del paradigma Web\textsuperscript{2} \cite{web2palazzi} ed esempi di applicazione. Opportunistic Sensing e applicazione dei principi al progetto. 

\section{Mobility Data Managment and Exploration}
Introduzione ai Mobility Data, cosa sono e a cosa possono servire. Panoramica sui requisiti necessari alla persistenza, all'elaborazione e all'interrogazione di questo tipo di dati. Dettaglio sulle tecniche di \emph{Trajectory Reconstruction}. Dettaglio sul tipo di problemi che risolve la gestione dei \emph{Mobility Data} con sistemi complessi tipo HERMES. Tutto materiale tratto da \cite{mdme} in particolare i primi capitoli e il cap. 3.
Perch� non abbiamo utilizzato questo approccio per lo sviluppo del progetto. 

\section{Versione precedente e lavori correlati}
Presentazione progetto Path 1.0 e valutazioni. Motivazioni sviluppo nuova architettura. 
Accenno ai progetti di Mobile Crowd Sensing (Crowd++ \cite{crowdplusplus} e Partecipact \cite{participact}).

\chapter{PathS Client} %------------------------------ PATHS CLIENT
\thispagestyle{empty}

Introduzione alla componente \emph{client} di \emph{PathS}. 

\section{Requisiti}
Presentazione degli obiettivi riguardanti la parte client:
\begin{itemize}
\item raccogliere campioni tramite sensore di luminosit� e microfono;
\item offrire indicazioni di percorso;
\item coinvolgere utente (tramite esperienza Augmented Reality).
\end{itemize}

\section{Tecnologie}
Riassunto delle tecnologie adottate per la realizzazione della componente \emph{client}: AndroidSDK, Wikitude.

\section{Stato dell'arte}
Risultati ottenuti con la parte client. Funzionamento ed eventuali limiti riscontrati.

\chapter{PathS Server} %------------------------------ PATHS SERVER
\thispagestyle{empty}

Introduzione alla componente \emph{server} di \emph{PathS}. 

\section{Requisiti}
Dettaglio dei requisiti per la parte server: raccolta campioni, persistenza, interfacciamento con il client, associazione dei campioni ai percorsi, servizio di routing.

\section{Componenti}
Identificazione delle componenti logiche del server sulla base delle funzioni che deve svolgere.
\subsection{Ricezione dei campioni}
La componente di ricezione campioni deve colloquiare con il client. Necessit� di definire un formato chiaro e \emph{standard}, facilmente implementabile con qualsiasi tecnologia \emph{client} ed eventualmente da altri sorgenti dati.
\subsection{Elaborazione dei campioni}
I campioni cos� come sono ricevuti dai client sono in formato grezzo e non consentono di sfruttare le inormazioni in esse contenute. Il passo principale per l'utilizzo dei dati � quello di associare ciascun dato ad un segmento della rete di trasporto. 
\subsection{Interrogazione}
Lo scopo pensato per l'utilizzo delle informazioni raccolte � quello di fornire servizi \emph{routing} alternativi, che tengano conto dei dati raccolti per suggerire percorsi che vanno oltre la semplice regola del percorso pi� breve. Il componente deve quindi:
\begin{itemize}
\item definire un formato di colloquio con il client per i percorsi;
\item implementare un algoritmo di routing;
\item modificare l'algoritmo di routing affinch� consideri le informazioni aggiuntive.
\end{itemize}

\section{Tecnologie e linee guida}
Tecnologie e librerie selezionate per l'applicazione globale al sistema server (Play! Framework, PostgresSQL, LeafletJS, Bootstrap). Presentazione dei criteri non funzionali con i quali � stato sviluppato il server (suddivisione delle responsabilit�, astrazione, estensibilit�).

\chapter{Ricezione dei campioni} %------------------------------ RICEZIONE CAMPIONI
\section{Formato}
Definizione e modifiche al formato di comunicazione.
\section{Implementazione}
Modalit� di implementazione e deserializzazione dei dati.
\section{Esempi}
Esempio di invio dal client e risultati ottenuti.

\chapter{Elaborazione dei campioni} %------------------------------ ELABORAZIONE CAMPIONI
\section{Servizi di Cartografia}
Requisiti per la selezione di un servizio di mappe su cui basarsi. Analisi comparativa e aspetti principali di:
\begin{itemize}
\item Google Maps
\item Openstreet Map
\end{itemize}
Il servizio di interrogazione \emph{Overpass API}, funzionalit� e tipo di \emph{query} implementata.

\section{Persistenza ed eleaborazione dei dati GIS}
Libreria scelta per la persistenza dei dati PostGIS, modello dati del DB e tipi utilizzati. Operazioni di interrogazione supportate dall'estensione del database.
Libreria scelta per la manipolazione dei dati in ambiente \emph{JAVA}:\emph{JTS}. Operazioni necessarie ed esempi di utilizzo. 

\section{Algoritmo di Map Matching}
Perch� � necessario eseguire una operazione di \emph{map matching}. A cosa servono questo tipo di algoritmi, panoramica sullo stato dell'arte e le soluzioni possibili.
\subsection{ST-MapMatching}
Esposizione delle caratteristiche principali dell'algoritmo come presentato in \cite{stmapmatching}.
\subsection{Pre-elaborazione del percorso - step 1}
Calcolo della bounding box, recupero della rete di trasporto coinvolta e inserimento a sistema.
\subsection{Identificazione dei candidati - step 2}
Definizione e modalit� di calcolo dei candidati. Operazioni implementate ed esempi di esecuzione.
\subsection{Valutazione e selezione dei candidati - step 3}
Applicazione della valutazione spazio-temporale dei campioni. Algoritmo di selezione dei candidati da assegnare ai campioni. 

\chapter{Routing} %------------------------------ ROUTING
Requisiti del servizio di routing. Definizione del formato di colloquio (GeoJSON) e interpretazione lato client.
\section{Implementazione}
Modalit� generali di implementazione e interfaccia \emph{Router}.
Pre-elaborazione e gestione dei segmenti a sistema tramite operazione di \emph{noding} e creazione della topologia (nodi e archi).
\section{Percorsi calcolati}
\subsection{Servizio Map Quest}
Interrogazione del servizio Map Quest come riscontro e \emph{fallback}. 
\subsection{Shortest Path}
Calcolo del percorso minimo con la libreria PGRouting.
\subsection{Percorsi con \emph{label}}
Calcolo dei percorsi valutando le label di luminosit� e rumorosit�. Viste e funzioni di supporto sul database. Modalit� di calcolo dei costi, spiegazione delle formule.

\chapter{Conclusione} %------------------------------ CONCLUSIONE
\section{Risultati}
Risultato complessivo del progetto. Funzionamento end-to-end con screenshot ed esempi.
\section{Miglioramenti ed Evoluzioni}
Limiti e compromessi dell'attuale implementazione. Aree di miglioramento ed evoluzione, tra cui:
\begin{itemize}
\item suddivisione ulteriore dei tratti di percorso;
\item alternative algoritmi di map-matching;
\item altri algoritmi di routing;
\item servizi aggiuntivi al routing (analisi dei posti pi� frequentati, funzioni ``get-together'' e \emph{social oriented}).
\end{itemize}

\backmatter

\begingroup %------------------------------ BIBLIOGRAPHY
  \makeatletter
  \let\ps@plain\ps@empty
  \makeatother
  \bibliography{tesi}
  \addcontentsline{toc}{chapter}{Bibliography}
  \bibliographystyle{ieeetr} % sort in order of appearance
\endgroup
\end{document} 